\documentclass[10pt,a4paper]{amsart}
\usepackage[margin=19mm]{geometry}
\usepackage{amsmath,amssymb,mathtools}
\usepackage[scr=rsfs,cal=euler]{mathalfa}
\usepackage{xfrac}
\usepackage[unicode,hidelinks,bookmarks=false]{hyperref}
\newcommand{\ie}{i.e.\ }
\newcommand{\cf}{cf.\ }
\newcommand{\resp}{respectively\ }
\newcommand{\Subsection}{\subsection}
\providecommand{\nobreakdash}{\nobreak\hbox{-}}
\setlength{\emergencystretch}{2em}
\multlinegap0pt
\usepackage{siunitx}
\usepackage{siunitx}
\usepackage{siunitx}
\usepackage{siunitx}
\usepackage{amssymb}
\usepackage{dsfont}
\usepackage{mathtools}
\usepackage{float}
\usepackage{siunitx}
\usepackage{tikz}
\usepackage{enumerate}
\usepackage{xcolor}
\newtheorem{lemma}{Lemma}
\DeclareMathOperator{\Gal}{Gal}
\newcommand{\Q}{\mathbb{Q}}
\newcommand{\Z}{\mathbb{Z}}
\newcommand{\fk}{\mathfrak }
\title{On the Iwasawa $\lambda$-invariant of the cyclotomic $\Z_2$-extension of a family of real quadratic fields in which $2$ splits}
\author{Josu\'e \'Avila, Foivos Chnaras}
\date{Source: arXiv:2605.09111v1 (CC BY 4.0)}
\begin{document}
\maketitle

\noindent\emph{Source and licence.} On the Iwasawa $\lambda$-invariant of the cyclotomic $\Z_2$-extension of a family of real quadratic fields in which $2$ splits. Version arXiv:2605.09111v1 (CC BY 4.0), distributed by arXiv under the Creative Commons Attribution 4.0 International licence, http://creativecommons.org/licenses/by/4.0/, as recorded in the arXiv licence metadata for this exact version and shown on the abstract page https://arxiv.org/abs/2605.09111v1. Full licence notice: \texttt{licenses/arxiv-2605.09111-CC-BY-4.0.txt}.\par\smallskip

\noindent\textit{Editorial scope.} Counted unit: the article's lemma P\_n (source0.tex lines 823--833). The article's complete original proof of it is delivered with it (source0.tex lines 835--1043, one \texttt{proof} environment of 209 lines). No mathematics is written, repaired or re-derived: the statement and the proof are the article's own lines, in the article's own notation, and no retained line is substituted or re-flowed.  No further source-local context is needed: every cross-reference of every retained line resolves inside the two delivered spans.  Nothing was omitted from any retained span, so the package carries no omission record.  No retained line contains a citation, so the wrapper carries no bibliography and no bibliography entry.  No retained line contains a figure environment, an image-inclusion command or drawing code, so no image is delivered and the package declares no asset dependency.  Editorial formatting only, in addition: an AMS \texttt{amsart} preamble that restates the article's own declarations and macros used by the retained lines, namely the environments \texttt{lemma} on separate counters, together with the standard packages those lines need.  The retained environments print in this document as Lemma 1 in order of appearance, whereas the article prints the same items with the numbering of its own full document.  The complete native source of the article as distributed in the arXiv e-print for this exact version is bundled unchanged as \texttt{sources/200195/source0.tex}.
\begin{lemma} \label{lem: property P_n}
Let $K$ be as above, and let $\epsilon_D$ be the fundamental unit of $K$.
Then
$$
(P_n)\ \forall n\ge 1
\iff
v_2(\log_2(\epsilon_D))=2
\iff
\epsilon_D\notin N_{K_1/K}(K_1^\times).
$$
\end{lemma}

\begin{proof}
Write
$$
2\mathcal O_K=\fk p_1\fk p_2.
$$
Since $2$ splits in $K$, we have
$$
K_{\fk p_1}\simeq K_{\fk p_2}\simeq \Q_2
$$ 
where $\Q_2$ represents the $2$-adic numbers.\\

Let $u_i\in \Q_2^\times$ be the image of $\epsilon_D$ in $K_{\fk p_i}$, and let
$\sigma$ be the non-trivial element of $\Gal(K/\mathbb Q)$. Since
$$
\sigma(\epsilon_D)=N_{K/\mathbb Q}(\epsilon_D)\,\epsilon_D^{-1},
$$
we get
$$
u_2=N_{K/\mathbb Q}(\epsilon_D)\,u_1^{-1}.
$$
As $\log_2(\pm 1)=0$, it follows that
$$
\log_2(u_2)=-\log_2(u_1),
$$
hence
$$
v_2(\log_2(u_1))=v_2(\log_2(u_2)).
$$
Thus the quantity $v_2(\log_2(\epsilon_D))$ is independent of the choice of dyadic place.

Set
$$
t:=v_2(\log_2(u_1))\in \Z_{\ge 2}\cup\{\infty\},
$$
where, as usual, $v_2(0)=\infty$.

For $u\in \Z_2^\times$, write
$$
u=\omega(u)\langle u\rangle,
\qquad
\omega(u)\in\{\pm1\},\quad \langle u\rangle\in 1+4\Z_2.
$$
Since
$$
\log_2:1+4\Z_2\longrightarrow 4\Z_2
$$
is an isomorphism of topological groups and $\log_2(x^{2^n})=2^n\log_2(x)$, we have
\begin{equation}\label{eq:2n-power-local}
u\in (\Z_2^\times)^{2^n}
\iff
\omega(u)=1\ \text{and}\ v_2(\log_2 u)\ge n+2.
\end{equation}

We first prove that
$$
(P_n)\ \forall n\ge 1
\iff
v_2(\log_2(\epsilon_D))=2.
$$



Assume first that $t=2$. Let
$$
x=(-1)^a\epsilon_D^m,
\qquad a\in\{0,1\},\quad 0\le m<2^n,
$$
represent a class in $\ker(\rho_n)$. Looking at the $\fk p_1$-component, the image of $x$ is
$$
(-1)^a u_1^m\in (\Z_2^\times)^{2^n}.
$$
By \eqref{eq:2n-power-local},
$$
v_2\!\bigl(\log_2((-1)^a u_1^m)\bigr)\ge n+2.
$$
Since $\log_2(-1)=0$, this becomes
$$
v_2(m\log_2(u_1))=v_2(m)+t\ge n+2.
$$
As $t=2$, we obtain
$$
v_2(m)\ge n,
$$
so $2^n\mid m$. Because $0\le m<2^n$, this forces $m=0$. Hence, $x=(-1)^a$.
But \eqref{eq:2n-power-local} also shows that $-1\notin (\Z_2^\times)^{2^n}$, so necessarily $a=0$.
Therefore $\ker(\rho_n)=0$; i.e. $(P_n)$ holds for every $n\ge 1$.

Conversely, suppose that $t>2$. If $t<\infty$, set
$$
n=t-1\ge 2;
$$
if $t=\infty$, choose any $n\ge 2$. Then the class of $\epsilon_D^2$ in
$$
E(K)/E(K)^{2^n}
$$
is non-trivial, since $\epsilon_D$ has infinite order and $2<2^n$. On the other hand, for each $i=1,2$ we have
$$
\omega(u_i^2)=1
$$
and
$$
v_2(\log_2(u_i^2))
=
v_2(2\log_2(u_i))
\ge n+2.
$$
Thus \eqref{eq:2n-power-local} gives
$$
u_i^2\in (\Z_2^\times)^{2^n}
\qquad (i=1,2).
$$
Hence, $\epsilon_D^2\in \ker(\rho_n)$, giving us a contradiction. Therefore, $t\le 2$.
Since always $t\ge 2$, we conclude that $t=2$.

This proves
$$
(P_n)\ \forall n\ge 1
\iff
v_2(\log_2(\epsilon_D))=2.
$$

We now prove that
$$
v_2(\log_2(\epsilon_D))=2
\iff
\epsilon_D\notin N_{K_1/K}(K_1^\times).
$$

Recall that
$$
\epsilon_D\in N_{K_1/K}(K_1^\times)
\iff
\epsilon_D\in N_{(K_1)_w/K_v}\bigl((K_1)_w^\times\bigr)
\quad\text{for every place }v\text{ of }K.
$$

At every finite place $v\nmid 2$, the extension $K_1/K$ is unramified or split, so every local unit is a norm; since $\epsilon_D$ is a unit, the local norm condition is automatic there. At the real places it is also automatic, because $2>0$. Hence only the dyadic places matter. Under the identifications
$$
K_{\fk p_i}\simeq \Q_2,
\qquad
(K_1)_{\fk P_i}\simeq \Q_2(\sqrt2),
$$
this becomes
$$
\epsilon_D\in N_{K_1/K}(K_1^\times)
\iff
u_i\in N_{\Q_2(\sqrt2)/\Q_2}\bigl(\Q_2(\sqrt2)^\times\bigr)
\qquad (i=1,2).
$$

By the local norm criterion for quadratic extensions,
$$
u_i\in N_{\Q_2(\sqrt2)/\Q_2}\bigl(\Q_2(\sqrt2)^\times\bigr)
\iff
(u_i,2)_2=1.
$$
Moreover,
$$
(u_2,2)_2
=
\bigl(N_{K/\Q}(\epsilon_D)u_1^{-1},2\bigr)_2
=
(u_1,2)_2,
$$
because $N_{K/\Q}(\epsilon_D)=\pm1$, $(\pm1,2)_2=1$, and $(u^{-1},2)_2=(u,2)_2$.
So it is enough to look at $u_1$.

For odd $u\in \Q_2^\times$, the explicit formula for the Hilbert symbol gives
$$
(u,2)_2=(-1)^{(u^2-1)/8}.
$$
Therefore,
$$
(u,2)_2=-1
\iff
u\equiv 3,5 \mod 8.
$$

On the other hand, if $u=\omega(u)\langle u\rangle$ with $\langle u\rangle\in 1+4\Z_2$, then
$$
\log_2(u)=\log_2(\langle u\rangle),
$$
and the series for $\log_2(1+x)$ with $x\in 4\Z_2$ shows that
$$
v_2(\log_2 u)=v_2(\langle u\rangle-1).
$$
Hence
$$
v_2(\log_2 u)=2
\iff
\langle u\rangle\equiv 5 \mod 8
\iff
u\equiv 3,5 \mod 8.
$$

Combining the last two equivalences, we obtain
$$
v_2(\log_2 u_1)=2
\iff
(u_1,2)_2=-1
\iff
u_1\notin N_{\Q_2(\sqrt2)/\Q_2}\bigl(\Q_2(\sqrt2)^\times\bigr).
$$
Since this is equivalent to
$$
\epsilon_D\notin N_{K_1/K}(K_1^\times),
$$
the proof is complete.
\end{proof}

\end{document}
